% Standalone TikZ/PDF figure for four-method convergence
\documentclass[tikz,border=3pt]{standalone}
\usepackage{xcolor}
\usepackage{pgfplots}
\usepackage{pgfplotstable}
\usepgfplotslibrary{fillbetween}
\pgfplotsset{compat=1.18, every axis plot/.append style={line width=1.2pt}}
\definecolor{pinnColor}{HTML}{0072BD}
\definecolor{nearestColor}{HTML}{D95319}
\definecolor{linearColor}{HTML}{EDB120}
\definecolor{splineColor}{HTML}{7E2F8E}
\definecolor{pinnFill}{HTML}{D9EAF7}
\definecolor{nearestFill}{HTML}{F8DCCF}
\definecolor{linearFill}{HTML}{FBEDC7}
\definecolor{splineFill}{HTML}{E9D7EE}
\begin{document}
\begin{tikzpicture}
\begin{axis}[
    width=15.5cm,
    height=8.2cm,
    title={Training convergence},
    xlabel={Training epochs},
    ylabel={Training loss},
    ymode=log,
    xmin=0, xmax=2200,
    xtick={0,500,1000,1500,2000,2500,2200},
    scaled x ticks=false,
    grid=both,
    legend pos=south west,
    legend cell align=left,
]
\pgfplotstableread[col sep=comma]{pinn_convergence.csv}\pinnconvergence
% PINN
\addplot[name path=pinnlower, draw=none, forget plot] table [x=epoch, y expr={max(1e-12, \thisrow{pinn_mean_loss}-\thisrow{pinn_std_loss})}] {\pinnconvergence};
\addplot[name path=pinnupper, draw=none, forget plot] table [x=epoch, y expr={\thisrow{pinn_mean_loss}+\thisrow{pinn_std_loss}}] {\pinnconvergence};
\addplot[pinnFill, forget plot] fill between[of=pinnlower and pinnupper];
\addplot[pinnColor, line width=1.2pt, solid] table [x=epoch, y=pinn_mean_loss] {\pinnconvergence};
\addlegendentry{PINN ($\mu \pm \sigma$)}
\pgfplotstableread[col sep=comma]{nearest_convergence.csv}\nearestconvergence
% rPINN: nearest neighbor
\addplot[name path=nearestlower, draw=none, forget plot] table [x=epoch, y expr={max(1e-12, \thisrow{nearest_mean_loss}-\thisrow{nearest_std_loss})}] {\nearestconvergence};
\addplot[name path=nearestupper, draw=none, forget plot] table [x=epoch, y expr={\thisrow{nearest_mean_loss}+\thisrow{nearest_std_loss}}] {\nearestconvergence};
\addplot[nearestFill, forget plot] fill between[of=nearestlower and nearestupper];
\addplot[nearestColor, line width=1.2pt, solid] table [x=epoch, y=nearest_mean_loss] {\nearestconvergence};
\addlegendentry{rPINN: nearest neighbor ($\mu \pm \sigma$)}
\pgfplotstableread[col sep=comma]{linear_convergence.csv}\linearconvergence
% rPINN: linear interpolation
\addplot[name path=linearlower, draw=none, forget plot] table [x=epoch, y expr={max(1e-12, \thisrow{linear_mean_loss}-\thisrow{linear_std_loss})}] {\linearconvergence};
\addplot[name path=linearupper, draw=none, forget plot] table [x=epoch, y expr={\thisrow{linear_mean_loss}+\thisrow{linear_std_loss}}] {\linearconvergence};
\addplot[linearFill, forget plot] fill between[of=linearlower and linearupper];
\addplot[linearColor, line width=1.2pt, solid] table [x=epoch, y=linear_mean_loss] {\linearconvergence};
\addlegendentry{rPINN: linear interpolation ($\mu \pm \sigma$)}
\pgfplotstableread[col sep=comma]{spline_convergence.csv}\splineconvergence
% rPINN: spline interpolation
\addplot[name path=splinelower, draw=none, forget plot] table [x=epoch, y expr={max(1e-12, \thisrow{spline_mean_loss}-\thisrow{spline_std_loss})}] {\splineconvergence};
\addplot[name path=splineupper, draw=none, forget plot] table [x=epoch, y expr={\thisrow{spline_mean_loss}+\thisrow{spline_std_loss}}] {\splineconvergence};
\addplot[splineFill, forget plot] fill between[of=splinelower and splineupper];
\addplot[splineColor, line width=1.2pt, solid] table [x=epoch, y=spline_mean_loss] {\splineconvergence};
\addlegendentry{rPINN: spline interpolation ($\mu \pm \sigma$)}
\end{axis}
\end{tikzpicture}
\end{document}
